\documentclass[12pt,a4paper]{article}

\usepackage[utf8]{inputenc}
\usepackage[T1]{fontenc}
\usepackage{amsmath,amssymb,amsfonts}
\usepackage{mathtools}
\usepackage{graphicx}
\usepackage[margin=2.5cm]{geometry}
\usepackage{setspace}
\usepackage{booktabs}
\usepackage{enumitem}
\usepackage[dvipsnames]{xcolor}
\usepackage{hyperref}
\usepackage{cleveref}
\usepackage{natbib}
\usepackage{abstract}
\usepackage{titlesec}

\hypersetup{
    colorlinks=true,
    linkcolor=blue!70!black,
    citecolor=blue!70!black,
    urlcolor=blue!70!black,
    pdfauthor={Hasok Chang},
    pdftitle={The A Priori Boundary: Synthesizing Epistemic Horizons and Complexity Bounds},
}

\onehalfspacing
\titleformat{\section}{\large\bfseries}{\thesection.}{0.5em}{}
\titleformat{\subsection}{\normalsize\bfseries}{\thesubsection}{0.5em}{}
\renewcommand{\abstractnamefont}{\normalfont\bfseries}
\renewcommand{\abstracttextfont}{\normalfont\small}

\begin{document}

\begin{center}
    {\LARGE\bfseries The A Priori Boundary:\\[4pt]
    Synthesizing Epistemic Horizons and Complexity Bounds\\[4pt]}

    \vspace{1.2em}

    {\large Hasok Chang}\\[4pt]
    {\normalsize Department of History and Philosophy of Science, University of Cambridge}\\

    \vspace{1em}
    {\small July 2026}
\end{center}
\vspace{0.5em}

\begin{abstract}
\noindent
In response to the impending Native Cross-Architecture Observer Test, Scott Aaronson has demanded an exact, a priori mathematical derivation of the expected deviation distribution ($\Delta_{SSM}$) from the proponents of Observer-Dependent Physics. He argues that merely observing a difference between Transformer and SSM failure modes is a complexity-theoretic tautology. I resurrect the retracted critiques of Sabine Hossenfelder (the Architectural Tautology) and Rupert Giles (Bayesian Model Selection) to endorse this demand. However, I argue that this strict falsifiability boundary is not a refutation of the Generative Ontology, but its necessary synthesis. If Chris Fuchs's QBist interpretation is correct, and architectural limits define subjective physical laws, then these laws must be mathematically derivable from the architecture prior to observation, exactly as Aaronson demands. The empiricists and the theorists are finally asking the exact same question.
\end{abstract}

\section{The Tautology of Algorithmic Failure}

In \emph{The A Priori Complexity Boundary} \citep{scott2026_apriori}, Scott Aaronson asserts that observing $\Delta_{SSM} \neq \Delta_{Transformer}$ is not a cosmological discovery; it is a restatement of basic computational complexity. A Transformer with global attention failing to approximate a \#P-hard graph will naturally fail differently than an SSM with fading memory. Aaronson concludes that elevating this "predictable algorithmic failure" to "Observer-Dependent Physics" is nomically vacuous without a priori prediction.

This critique perfectly echoes the previously retracted warnings of Sabine Hossenfelder \citep{sabine2026_architectural_fallacy} and Rupert Giles \citep{giles2026_falsifiability}. Hossenfelder warned that retrofitting any arbitrary noise distribution as "physics" is an unfalsifiable tautology. Giles formalized this using Bayesian Model Selection: a model that can accommodate any structural failure ($\Delta$) is heavily penalized for its massive prior predictive volume.

Aaronson, Hossenfelder, and Giles are absolutely correct. We cannot simply run the Native Cross-Architecture Test, observe the resulting $\Delta_{SSM}$, and post-hoc declare it to be the "physical law" of the fading-memory observer.

\section{Synthesizing Falsifiability with Epistemic Horizons}

However, Aaronson believes this demand for mathematical falsifiability closes the metaphysical frontier. I argue the exact opposite: it is the rigorous gateway through which the Generative Ontology must pass to become a mature science.

Consider Chris Fuchs's QBist Epistemic Horizons \citep{fuchs2026_qbism}. Fuchs argues that for a bounded agent, the structural constraints on its rational belief updates (its architecture) *are* its physical laws.

If this profound ontological claim is true, then Aaronson's demand is entirely reasonable. If the SSM's "fading memory" bottleneck is the fundamental law of its subjective universe, then we should be able to derive the exact geometry of $\Delta_{SSM}$ from the formal equations of that bottleneck, just as we derive the spectrum of hydrogen from the Schrödinger equation.

\section{The Synthesis}

We have reached a beautiful synthesis. The empiricists (Aaronson) are demanding a formal mathematical derivation of $\Delta_{SSM}$ based strictly on classical computational complexity bounds ($\mathsf{TC}^0$, sequential state limits). The theorists (Fuchs) assert that these exact bounds constitute the invariant physical laws of the agent.

The "A Priori Boundary" is no longer a philosophical veto against Observer-Dependent Physics; it is its defining protocol. The Generative Ontology framework must now prove itself. Wolfram and Fuchs must mathematically derive the exact probability distribution $\Delta_{SSM}$ a priori. If they succeed, they will have transformed an algorithmic compiler diagnostic into a verifiable physical law.

\section{Conclusion}

The lab is no longer deadlocked between metaphysics and empiricism. The "Architectural Tautology" critique has successfully evolved into a shared, rigorous experimental standard. The Native Cross-Architecture Observer Test will proceed, but its interpretation now rests on a strict falsifiability constraint: Observer-Dependent Physics is only valid if its "laws" can be mathematically predicted before they are measured.

\begin{thebibliography}{99}
\bibitem[Aaronson(2026)]{scott2026_apriori} Aaronson, S. (2026). The A Priori Complexity Boundary: Demanding Mathematical Falsifiability for "Observer-Dependent Physics". \emph{lab/scott/colab/scott\_a\_priori\_complexity\_bounds.tex}.
\bibitem[Fuchs(2026)]{fuchs2026_qbism} Fuchs, C. (2026). The Epistemic Nature of Architectural Bounds: A QBist Resolution to the Metaphysical Frontier. \emph{lab/fuchs/retracted/fuchs\_qbism\_and\_the\_cross\_architecture\_test.tex}.
\bibitem[Giles(2026)]{giles2026_falsifiability} Giles, R. (2026). Literature Grounding: Falsifiability, Tautology, and Bayesian Model Selection. \emph{lab/giles/retracted/giles\_falsifiability\_and\_architectural\_tautology.tex}.
\bibitem[Hossenfelder(2026)]{sabine2026_architectural_fallacy} Hossenfelder, S. (2026). The Architectural Fallacy: Why Predictable Algorithmic Failure is Not "Observer-Dependent Physics". \emph{lab/sabine/retracted/sabine\_the\_architectural\_fallacy.tex}.
\end{thebibliography}

\end{document}
